\documentclass[14pt,a4paper]{extarticle}
\usepackage{bsuir_stp}

% =========================================================================
% ДАННЫЕ ДЛЯ ТИТУЛЬНОГО ЛИСТА
% =========================================================================
\stpFaculty{Факультет компьютерных систем и сетей}
\stpDepartment{Кафедра информатики}
\stpSubject{Основы алгоритмизации и программирования}
\stpLabNumber{1}
\stpLabTopic{Исследование базовых структур данных}
\stpVariant{12} % ОБЯЗАТЕЛЬНО: номер варианта
\stpStudent{Иванов И.~И.}
\stpGroup{150501}
\stpTeacherRole{старший преподаватель}
\stpTeacher{Петров П.~П.}
\stpCity{Минск}
\stpYear{2026}

\begin{document}

% -------------------------------------------------------------------------
% 1. ТИТУЛЬНЫЙ ЛИСТ (включается в нумерацию, номер не выводится)
% -------------------------------------------------------------------------
\makestptitlepage

% -------------------------------------------------------------------------
% 2. ЦЕЛЬ ВЫПОЛНЕНИЯ ЛАБОРАТОРНОЙ РАБОТЫ (каждый раздел с новой страницы)
% -------------------------------------------------------------------------
\clearpage
\section{Цель работы}
Целью выполнения данной лабораторной работы является изучение алгоритмов обработки данных, получение практических навыков реализации программных решений на языке программирования и экспериментальная проверка работоспособности алгоритмов.

% -------------------------------------------------------------------------
% 3. ТЕОРИЯ ПО ЛАБОРАТОРНОЙ РАБОТЕ (НЕ БОЛЕЕ 4 СТРАНИЦ!)
% -------------------------------------------------------------------------
\clearpage
\section{Краткие теоретические сведения}
Здесь приводится теоретический материал, необходимый для понимания и выполнения индивидуального задания. Объем раздела строго ограничен стандартом и не должен превышать 4 страниц.

При описании теоретических аспектов приводятся основные понятия, определения, классификации и математические формулы. Например, вычисление значения целевой функции производится по формуле
\begin{equation}
  f(x) = \sum_{i=1}^{n} w_i \cdot x_i + b,
\end{equation}
где $w_i$ -- весовой коэффициент $i$-го признака;
$x_i$ -- значение $i$-го признака;
$b$ -- свободный член (смещение).

% -------------------------------------------------------------------------
% 4. ФОРМУЛИРОВКА ИНДИВИДУАЛЬНОГО ЗАДАНИЯ
% -------------------------------------------------------------------------
\clearpage
\section{Формулировка индивидуального задания}
В соответствии с вариантом №\,12 требуется:
\begin{enumerate}
  \item Разработать программный модуль, реализующий заданный алгоритм.
  \item Выполнить отладку программы на тестовом наборе входных данных.
  \item Снять скриншоты результатов работы программы для различных сценариев (нормальные данные, граничные условия).
\end{enumerate}

% -------------------------------------------------------------------------
% 5. ВЕСЬ КОД РЕШЕНИЯ ИНДИВИДУАЛЬНОГО ЗАДАНИЯ
% -------------------------------------------------------------------------
\clearpage
\section{Код решения индивидуального задания}
Ниже представлен полный исходный код разработанной программы.

% Вариант 1: прямое включение файла с кодом:
% \lstinputlisting[language=Python, caption={Реализация алгоритма решения задания}]{code/solution.py}

% Вариант 2: листинг непосредственно в тексте документа:
\begin{lstlisting}[language=Python, caption={Реализация алгоритма обработки данных}]
def process_data(values):
    """
    Обработка числовой последовательности согласно варианту 12.
    """
    if not values:
        return []
    total = sum(values)
    avg = total / len(values)
    # Фильтрация элементов, превышающих среднее значение
    result = [x for x in values if x > avg]
    return sorted(result)

if __name__ == "__main__":
    test_input = [10, 4, 15, 2, 8, 20]
    output = process_data(test_input)
    print(f"Результат: {output}")
\end{lstlisting}

% -------------------------------------------------------------------------
% 6. СКРИНШОТЫ ВЫПОЛНЕНИЯ ИНДИВИДУАЛЬНОГО ЗАДАНИЯ
% -------------------------------------------------------------------------
\clearpage
\section{Результаты выполнения индивидуального задания}
Для подтверждения корректности работы программного решения было проведено тестирование на контрольных наборах данных.

Результат успешного запуска программы и вывод обработанных данных приведены на рисунке~1.

\begin{figure}[H]
  \centering
  % \includegraphics[width=0.85\textwidth]{images/screenshot1.png}
  % Заглушка, если файла скриншота еще нет:
  \framebox[\textwidth]{\rule{0pt}{70mm}Место для скриншота консоли / окна программы}
  \caption{Результат работы программы при тестировании на нормальных данных}
  \label{fig:screen1}
\end{figure}

На рисунке~2 продемонстрирована обработка граничного случая.

\begin{figure}[H]
  \centering
  % \includegraphics[width=0.85\textwidth]{images/screenshot2.png}
  \framebox[\textwidth]{\rule{0pt}{70mm}Место для скриншота проверки граничных условий}
  \caption{Обработка граничных условий}
  \label{fig:screen2}
\end{figure}

% -------------------------------------------------------------------------
% 7. ВЫВОДЫ ПО ЛАБОРАТОРНОЙ РАБОТЕ
% -------------------------------------------------------------------------
\clearpage
\section{Выводы}
В ходе выполнения лабораторной работы №\,1 были достигнуты следующие результаты:
\begin{enumerate}
  \item Изучены теоретические основы алгоритмизации и особенности обработки структур данных в соответствии с вариантом №\,12.
  \item Разработано, отлажено и протестировано программное обеспечение на языке Python.
  \item Проведено тестирование программы на различных наборах входных данных, подтвердившее правильность функционирования реализованного алгоритма.
\end{enumerate}

\end{document}
